\section{Conclusion}\label{sec:conclusion}

Households exposed to large pork-price shocks systematically
overpredict future inflation.
The forecast-error coefficient ranges from $-5.62$ under
province-plus-wave fixed effects to $-7.40$ under the most
demanding province-plus-region-by-wave specification, exceeding
the rational benchmark ($|\beta_3^R| \leq 0.05$) by a wide
margin and surviving wild cluster bootstrap inference with
31 province clusters.
A formal Wald test rejects equality of meat and grain
forecast-error coefficients ($p = 0.071$), confirming that
the pattern is specific to pork, not a generic food-price
confound.

The paper makes three contributions.
First, the forecast-error test uses an external commodity shock
measured by the NBS rather than within-survey belief revisions.
This eliminates the mechanical negative correlation between
revisions and forecast errors that contaminates standard
\citet{Coibion2015}-style designs, allowing the sign of
$\hat{\beta}_3$ to be interpreted as evidence of overreaction
without caveat.

Second, the signal-extraction benchmark provides a formal
separation between rational forecast errors due to
transitory-shock uncertainty and behavioural overreaction.
The rational coefficient is bounded by $|\beta_3^R| \leq
\omega \approx 0.05$, regardless of beliefs about persistence.
The estimated coefficient exceeds this bound by a large factor
under all specifications, though the ordinal expectation scale
inflates the exact ratio and the point estimate of $\theta$
should be read as an order-of-magnitude statement.

Third, the commodity placebo test disciplines the mechanism
as salience-specific.
Grain shocks of comparable CPI weight do not produce the same
expectation response, isolating salience---driven by purchase
frequency and media coverage---as the operative channel.

The results also extend the diagnostic expectations literature.
\citet{Bordalo2018} develop the theoretical framework;
\citet{Bordalo2020b} provide evidence from professional
forecasters.
The household-level evidence here shows that diagnostic
distortions are not confined to professionals and pins down
the source: a salient signal about a narrow price category
is extrapolated to the aggregate price level.

Several limitations apply.
The CFPS is biennial, so the analysis cannot trace within-year
adjustment dynamics.
Province-level meat CPI conflates supply and demand movements;
ASF provides a plausible supply instrument, but the treatment
is an equilibrium price.
The ordinal expectation measure limits the precision of the
expectation-channel estimate, and with 31 province clusters
all inference relies on the wild cluster bootstrap.

Two extensions follow naturally.
First, higher-frequency household expectation data matched
to commodity-price signals at finer geographic resolution
would sharpen identification of adjustment dynamics and
allow estimation of the salience effect's decay rate.
The Michigan Survey of Consumers could support a parallel
design for the United States using gasoline prices.
Second, a structural model of belief formation under diagnostic
expectations, calibrated to the reduced-form estimates here,
could quantify the welfare cost of the expectation wedge and
evaluate whether targeted communication about relative-price
versus aggregate-price movements can reduce overreaction.
